% ============================================================================
% Fichier  : partie1/P01_C04_fraude.tex
% Titre    : La Fraude dans et par les Systèmes d'Information
% Auteur   : MINKA MI NGUIDJOÏ Thierry Emmanuel
% Version  : 1.0 -- octobre 2026
% ============================================================================
% Sources : COSO et ACFE, Guide de gestion du risque de fraude, synthèse
% (décembre 2016, dossier Normes) ; Cressey (1953) pour le triangle de la
% fraude ; COSO 2013 (principe 8) ; article A0 (axiome 4, interaction
% stratégique) et article A1 (usage stratégique de NEK). L'obligation de
% révélation du commissaire aux comptes (AUSCGIE) est décrite sans numéro
% d'article : à vérifier dans le texte.
% ============================================================================

\chapter{La Fraude dans et par les Systèmes d'Information}
\label{chap:fraude}
\chaptermark{La fraude et les SI}

\epigraph{%
    « L'erreur se cache par hasard.\\
    La fraude se cache exprès~: c'est ce qui la rend auditable. »%
}{Principe de l'auditeur}

\niveauChapitre{Intermédiaire}{%
    Maîtrise des Chapitres~\ref{chap:controle-interne} (séparation des
    tâches) et~\ref{chap:risque} (cartographie des risques)%
}

\begin{boxobjectifs}
\begin{itemize}
    \item Définir la fraude et la distinguer de l'erreur.
    \item Expliquer le triangle de la fraude et le relier à l'interaction
          stratégique entre auditeur et audité.
    \item Classer les schémas de fraude assistés par l'informatique et leurs
          signaux dans les données.
    \item Concevoir un dispositif de gestion du risque de fraude selon le
          guide du COSO et de l'ACFE.
    \item Reconnaître un usage stratégique de la réponse \NEK{}.
    \item Savoir quand l'audit doit s'arrêter et laisser place à
          l'enquête.
\end{itemize}
\end{boxobjectifs}

\bigskip

% ============================================================================
\section*{Question directrice : erreur ou fraude~?}
% ============================================================================

Au Chapitre~\ref{chap:risque}, trois risques intolérables de LiZbiZ
supposaient une intention~: ajuster le stock, payer des fournisseurs
fictifs, utiliser le compte d'un agent parti. Lorsque AUDIX CONSEIL
présente sa cartographie, le responsable informatique proteste~: « Ce sont
des erreurs de saisie, pas des fraudes. Nos agents sont honnêtes. »

Peut-être. Mais l'auditeur ne présume ni l'honnêteté ni la malhonnêteté~: il
cherche des faits. Ce chapitre définit la fraude
(section~\ref{sec:c04-definir}), explique pourquoi elle survient
(section~\ref{sec:c04-triangle}), décrit ses schémas dans les systèmes
d'information (section~\ref{sec:c04-schemas}), présente le dispositif qui la
prévient et la détecte (section~\ref{sec:c04-dispositif}), montre comment
la repérer dans les données et dans les réponses
(sections~\ref{sec:c04-donnees} et~\ref{sec:c04-strategie}), et fixe la
limite entre audit et enquête (section~\ref{sec:c04-enquete}).

\separateur

% ============================================================================
\section{Définir la Fraude}
\label{sec:c04-definir}
% ============================================================================

\begin{boxdef}
\termecle{Fraude}\index{fraude}~: acte \emph{intentionnel}, commis par une ou
plusieurs personnes de l'organisation ou extérieures à elle, qui recourt à
la tromperie pour obtenir un avantage indu ou causer un préjudice.

\medskip
\termecle{Erreur}\index{erreur}~: inexactitude \emph{non intentionnelle}
(faute de saisie, mauvaise interprétation d'une règle, panne).
\end{boxdef}

Ce qui distingue la fraude de l'erreur, c'est l'intention, et donc la
\emph{dissimulation}. Une erreur laisse des traces au hasard~; une fraude
efface les siennes, ou les place là où personne ne regarde. C'est pourquoi
la fraude obéit à des schémas, que l'auditeur peut connaître.

\begin{tableaucours}{L{3.6cm}L{4.4cm}L{4.2cm}}{%
    \tableauHeader{Famille} &
    \tableauHeader{Définition} &
    \tableauHeader{Exemple chez LiZbiZ}
}
Détournement d'actifs & Appropriation de biens ou de fonds de
l'organisation & Vol de marchandises masqué par un ajustement de stock \\
Corruption & Usage de sa fonction pour obtenir un avantage indu, souvent
avec un tiers & Commission versée par un transporteur pour être
retenu \\
Fraude à l'information financière & Présentation trompeuse des états
financiers & Stock surévalué pour améliorer le résultat \\
\end{tableaucours}
\vspace{-0.8em}
\begin{center}{\small\itshape Tableau~4.1. Trois familles de fraude
(classification usuelle de la profession)}\end{center}

% ============================================================================
\section{Pourquoi la Fraude Survient}
\label{sec:c04-triangle}
% ============================================================================

\subsection{Le triangle de la fraude}

Dans les années 1950, le criminologue Donald Cressey étudie des personnes
condamnées pour détournement de fonds~\cite{cressey1953}. Ses travaux ont
donné naissance au \termecle{triangle de la fraude}\index{triangle de la
fraude}~: la fraude survient lorsque trois conditions sont réunies.

\noindent\begin{minipage}{\linewidth}
\begin{center}
\begin{tikzpicture}[
    sommet/.style = {
        draw=CouleurPrimaire, fill=CouleurDef, rounded corners=4pt,
        text width=3.2cm, align=center, minimum height=1.0cm,
        font=\small\bfseries\color{CouleurPrimaire}
    }
]
\coordinate (A) at (0,3.0);
\coordinate (B) at (-3.0,0);
\coordinate (C) at (3.0,0);
\draw[very thick, color=CouleurAccent] (A) -- (B) -- (C) -- cycle;
\node[sommet] at (A) {Pression\\ \mdseries\footnotesize dette, objectifs intenables};
\node[sommet] at (B) {Opportunité\\ \mdseries\footnotesize contrôle faible, droits cumulés};
\node[sommet] at (C) {Rationalisation\\ \mdseries\footnotesize « je le rendrai », « on me doit »};
\node[font=\footnotesize\itshape, color=CouleurSecondaire] at (0,1.0) {Fraude};
\end{tikzpicture}
\end{center}
\vspace{-0.4em}
\begin{center}{\small\itshape Figure~4.1. Le triangle de la fraude, d'après
les travaux de Cressey}\end{center}
\end{minipage}

\medskip

L'auditeur et le contrôleur interne n'ont guère de prise sur la pression ni
sur la rationalisation. Ils agissent sur l'\textbf{opportunité}~: c'est
pourquoi la séparation des tâches et la revue des droits
(Chapitre~\ref{chap:controle-interne}) sont les premières défenses contre la
fraude. Des auteurs ont ajouté un quatrième élément, la \emph{capacité}~: la
position et les compétences qui permettent de passer à l'acte. Dans les
systèmes d'information, cette capacité est souvent technique~: un
administrateur peut faire ce que le contrôle interne interdit aux autres.

\subsection{Un jeu entre l'auditeur et l'audité}

L'auditologie pose que l'auditeur et l'audité poursuivent des intérêts
distincts, et que leurs actions se comprennent comme des coups dans un jeu
à information imparfaite (axiome~4~\cite{minka2026auditologie}). La fraude en
est le cas extrême~: le fraudeur sait ce que l'auditeur ignore, et cherche à
ce qu'il continue de l'ignorer. Deux conséquences pratiques en découlent.

\begin{enumerate}
    \item \textbf{L'audité adapte ses réponses à l'enjeu.} Une même personne
          ne répond pas de la même façon dans une auto-évaluation et dans un
          entretien formel (section~\ref{sec:c04-strategie}).
    \item \textbf{L'auditeur doit être imprévisible.} Un programme de tests
          connu d'avance indique au fraudeur où ne pas laisser de traces~:
          les échantillons, les dates et certains tests doivent rester
          inconnus de l'audité.
\end{enumerate}

% ============================================================================
\section{Les Schémas de Fraude Assistés par l'Informatique}
\label{sec:c04-schemas}
% ============================================================================

\begin{tableaucours}{L{2.9cm}L{3.2cm}L{3.0cm}L{3.1cm}}{%
    \tableauHeader{Schéma} &
    \tableauHeader{Vulnérabilité exploitée} &
    \tableauHeader{Contrôle préventif} &
    \tableauHeader{Signal dans les données}
}
Fournisseur fictif & Cumul création du fournisseur et paiement & Séparation
des tâches, validation des nouveaux fournisseurs & Coordonnées bancaires
partagées avec un salarié, absence d'adresse \\
Modification des coordonnées bancaires & Données de référence modifiables
sans validation & Double validation, journal des modifications &
Changement de RIB suivi d'un paiement dans les jours suivants \\
Fausse réception & Cumul commande et réception & Réception confirmée par
une personne distincte & Réception sans bon de transport, à l'heure de
fermeture \\
Ajustement de stock & Cumul réception et ajustement & Validation au-delà
d'un seuil & Ajustements concentrés en fin de mois, sans pièce \\
Fractionnement & Seuil de validation connu & Contrôle des cumuls par
fournisseur & Paiements juste sous le seuil, répétés \\
Compte dormant & Comptes d'agents partis non désactivés & Revue périodique
des comptes & Connexion après la date de départ \\
Abus d'administrateur & Droits techniques étendus, sans trace & Comptes
nominatifs, journalisation protégée & Opérations de gestion faites par un
compte d'administration \\
\end{tableaucours}
\vspace{-0.8em}
\begin{center}{\small\itshape Tableau~4.2. Schémas de fraude et
signaux}\end{center}

\begin{boxalert}
\textbf{Violation typique.} Laisser le responsable informatique, qui détient
tous les droits, administrer lui-même les journaux qui tracent ses propres
actions.

\smallskip
\textbf{Règle~:} les journaux des comptes d'administration sont conservés
hors de leur portée et examinés par une personne indépendante.
\end{boxalert}

% ============================================================================
\section{Le Dispositif Anti-Fraude}
\label{sec:c04-dispositif}
% ============================================================================

Le principe~8 du COSO~2013 demande de prendre en compte le risque de fraude
(Chapitre~\ref{chap:controle-interne}). Le COSO et l'Association of
Certified Fraud Examiners (\sigle{ACFE}) ont publié en 2016 un guide qui
décline cette exigence en un programme global de gestion du risque de
fraude~\cite{coso2016fraude}. Ses cinq principes sont alignés sur les cinq
composantes du contrôle interne.

\begin{tableaucours}{L{3.6cm}L{8.6cm}}{%
    \tableauHeader{Volet} &
    \tableauHeader{Contenu (reformulé)}
}
Gouvernance & Une politique écrite de gestion du risque de fraude, portée
par l'organe de gouvernance et la direction \\
Évaluation & Une évaluation régulière et complète du risque de fraude,
schéma par schéma \\
Activités de contrôle & Des contrôles préventifs et détectifs choisis pour
les schémas identifiés \\
Investigation et correction & Un dispositif de signalement et une approche
coordonnée des investigations et des mesures correctives \\
Pilotage & Des évaluations qui vérifient que les cinq volets fonctionnent,
et la communication des déficiences \\
\end{tableaucours}
\vspace{-0.8em}
\begin{center}{\small\itshape Tableau~4.3. Les cinq volets de la gestion du
risque de fraude (d'après COSO et ACFE, 2016)}\end{center}

\begin{boxmetier}{Contrôleur interne}
\textbf{Bâtir le programme anti-fraude de LiZbiZ}
\begin{enumerate}
    \item Faire adopter par le conseil une politique anti-fraude~: tolérance
          zéro, obligation de signaler, protection des lanceurs d'alerte.
    \item Dresser, schéma par schéma, la liste des fraudes possibles dans
          l'application intégrée (Tableau~4.2).
    \item Supprimer d'abord les opportunités~: cumuls de droits, comptes
          dormants, seuils connus et non surveillés.
    \item Mettre en place des contrôles détectifs automatiques sur les
          données, revus chaque mois.
    \item Ouvrir un canal de signalement confidentiel, hors de la
          hiérarchie directe.
\end{enumerate}
\end{boxmetier}

% ============================================================================
\section{Repérer la Fraude dans les Données}
\label{sec:c04-donnees}
% ============================================================================

L'informatique facilite la fraude, mais elle en garde la trace. L'auditeur
analyse l'ensemble des opérations, et non un échantillon, à la recherche des
signaux du Tableau~4.2~: doublons, valeurs juste sous un seuil, opérations
hors des heures de travail, coordonnées partagées entre un fournisseur et un
salarié. Le Chapitre~\ref{chap:audit-continu} présentera les outils de ces
analyses.

\begin{boxscenario}{Opération LIZBIZ. Ce que disent les données}
\textbf{Tests réalisés} par AUDIX CONSEIL sur les extractions de l'application
intégrée, pour 2025 et le premier trimestre~2026.

\medskip
\textbf{Constatation 1.} Le fournisseur « TRANSLOG SERVICES », créé en
mars~2025 par Mme~OWONA, a reçu 23~paiements de 4,9~millions FCFA chacun,
soit 112,7~millions FCFA. Le seuil de double validation est de 5~millions.
Ses coordonnées bancaires sont identiques à celles qu'un salarié a déclarées
pour sa paie.

\medskip
\textbf{Constatation 2.} À l'entrepôt de Douala, les ajustements négatifs de
stock saisis par M.~FOUDA s'élèvent à 148~millions FCFA, dont 71~\% le
dernier jour ouvré du mois. Aucun n'est appuyé par un procès-verbal
d'inventaire.

\medskip
\textbf{Constatation 3.} Sur l'écart de stock de 186~millions FCFA, 38~millions
ne peuvent être rattachés à aucune opération~: les extractions des autres
entrepôts sont incomplètes.

\medskip
\textbf{Analyse} (confiance moyenne). Les constatations 1 et 2 réunissent
les signaux de deux schémas~: fournisseur fictif avec fractionnement, et
ajustement de stock. Elles ne prouvent pas une fraude, car l'intention reste
à établir, mais elles imposent d'en informer le commanditaire. La
constatation 3 est cotée \NEK{}~: on ne sait pas, et il faut dire pourquoi.
\end{boxscenario}

% ============================================================================
\section{Repérer la Dissimulation dans les Réponses}
\label{sec:c04-strategie}
% ============================================================================

Le prologue l'a annoncé~: les réponses données en entretien formel ont
différé de celles de l'auto-évaluation. La réponse \NEK{} peut traduire une
ignorance sincère~; elle peut aussi servir à ne pas répondre. L'article
\emph{Modelling Epistemic Uncertainty in IS Audit} propose un indice, le
\termecle{SUI}\index{SUI (indice d'usage stratégique)} (\textit{Strategic Use
Index}), qui compare la distribution des \NEK{} entre questions à fort enjeu
et à faible enjeu, pour distinguer l'ignorance authentique de la
dissimulation~\cite{minka2026nek}. Le Chapitre~\ref{chap:incertitude} en
donnera la définition.

\begin{boxlecture}{Entrepôt de Douala}
Les mêmes dix questions sur la gestion du stock de Douala sont posées aux
mêmes deux responsables, d'abord en auto-évaluation écrite, puis en
entretien formel, après la découverte des ajustements.

\medskip
\begin{tableaucours}{L{3.6cm}ccccc}{%
    \tableauHeader{Contexte} &
    \tableauHeader{\OUI{}} & \tableauHeader{\NON{}} & \tableauHeader{\NEK{}} &
    \tableauHeader{$U$} & \tableauHeader{$H$ (bit)}
}
Auto-évaluation & 17 & 2 & 1 & 0,05 & 0,75 \\
Entretien formel & 9 & 3 & 8 & 0,40 & 1,46 \\
\end{tableaucours}

\medskip
\textbf{Constatation.} Le taux de \NEK{} passe de 5~\% à 40~\% lorsque
l'enjeu augmente, pour les mêmes personnes et les mêmes questions.

\medskip
\textbf{Analyse} (confiance faible à moyenne). Une ignorance sincère ne
devrait pas dépendre du contexte. Ce déplacement est un \emph{signal}
d'usage stratégique de \NEK{}, pas une preuve~: d'autres explications sont
possibles, comme la crainte de se tromper devant un auditeur. Il oriente
les tests vers les questions où \NEK{} a le plus augmenté.

\medskip
\hfill\lecture{0{,}45}{0{,}40}{1{,}46}
\end{boxlecture}

% ============================================================================
\section{De l'Audit à l'Enquête}
\label{sec:c04-enquete}
% ============================================================================

L'auditeur n'est ni un enquêteur ni un juge. Lorsque ses travaux révèlent des
indices sérieux de fraude, il ne poursuit pas seul.

\begin{boxprocedure}
\textbf{Conduite à tenir devant des indices de fraude.} Le papier de travail
de l'auditeur doit permettre d'établir qu'il a~:
\begin{itemize}
    \item conservé les extractions et les journaux utilisés, avec leur
          empreinte numérique, sans les modifier~;
    \item consigné les constatations séparément de leur analyse, sans
          qualification pénale («~fraude~», «~vol~»)~;
    \item informé le commanditaire désigné par la lettre de mission, par
          écrit, et non les personnes mises en cause~;
    \item suspendu les tests susceptibles d'alerter les personnes en cause,
          dans l'attente de la décision du commanditaire.
\end{itemize}
\end{boxprocedure}

La suite relève d'une enquête interne, d'un expert judiciaire ou de la
justice. Elle obéit aux règles de la preuve numérique~: chaîne de
conservation, intégrité, opposabilité. Le Chapitre~\ref{chap:audit-ia}
et le cours de Théorie et Pratique de l'Investigation Numérique y
reviennent.

\begin{boxmetier}{Commissaire aux comptes}
\textbf{Des indices de fraude dans l'audit légal}

Le commissaire aux comptes de LiZbiZ, informé des constatations d'AUDIX
CONSEIL, doit en apprécier l'effet sur les comptes~: stock surévalué,
charges fictives. Le droit OHADA des sociétés commerciales lui impose en
outre de révéler au ministère public les faits délictueux dont il a
connaissance dans l'exercice de sa mission~\cite{ohada2014auscgie}
(l'article est à vérifier dans le texte avant toute citation). Il ne peut
donc pas «~régler l'affaire en interne~».
\end{boxmetier}

\begin{boiteRemarque}{L'auditeur doit-il chercher la fraude~?}
Pas systématiquement~: la lettre de mission d'AUDIX CONSEIL l'exclut
(Chapitre~\ref{chap:audit-auditologie}). Mais il doit tenir compte du risque
de fraude dans la conception de ses tests, rester attentif aux signaux, et
ne jamais ignorer un indice qu'il a rencontré. Ne pas chercher n'est pas
fermer les yeux.
\end{boiteRemarque}

\separateur

\begin{boxresume}
\textbf{Ce qu'il faut retenir de ce chapitre~:}
\begin{itemize}
    \item \textbf{La fraude est intentionnelle et se dissimule}~; c'est ce
          qui la distingue de l'erreur et ce qui lui donne des schémas
          reconnaissables.
    \item \textbf{Le triangle de la fraude}~: pression, opportunité,
          rationalisation~; l'auditeur et le contrôleur agissent surtout sur
          l'opportunité.
    \item \textbf{Les schémas assistés par l'informatique} exploitent des
          cumuls de droits, des données de référence non protégées, des seuils
          connus et des comptes dormants~; chacun laisse un signal dans les
          données.
    \item \textbf{Un programme anti-fraude} a cinq volets~: gouvernance,
          évaluation, contrôles, investigation et correction, pilotage.
    \item \textbf{La réponse \NEK{}} peut être stratégique~: un taux de \NEK{}
          qui dépend de l'enjeu est un signal, pas une preuve.
    \item \textbf{Devant des indices}, l'auditeur conserve les preuves,
          informe le commanditaire et laisse la place à l'enquête.
\end{itemize}
\end{boxresume}

\bigskip

\begin{boxexo}[title={\ding{45}\quad Exercice~4.1. Erreur ou fraude~? \niveaubase}]
Pour chacune des situations suivantes, indiquez s'il s'agit plutôt d'une
erreur ou d'un signal de fraude, et le test qui permettrait de trancher.
\begin{enumerate}[(a)]
    \item Une facture est saisie deux fois, avec deux numéros différents.
    \item Les coordonnées bancaires d'un transporteur sont modifiées, puis
          rétablies trois jours plus tard, après un paiement.
    \item Un ajustement de stock positif compense exactement un ajustement
          négatif de la veille.
    \item Une commande est passée un dimanche par un compte d'administration.
\end{enumerate}
\end{boxexo}

\begin{boxexo}[title={\ding{45}\quad Exercice~4.2. Le triangle de TRANSLOG SERVICES \niveaumoyen}]
Reprenez la constatation 1 de la section~\ref{sec:c04-donnees}.
\begin{enumerate}[(a)]
    \item Identifiez l'élément du triangle de la fraude sur lequel LiZbiZ peut
          agir, et les contrôles correspondants.
    \item Rédigez la constatation de façon factuelle, sans qualification
          pénale.
    \item Proposez deux tests complémentaires qui confirmeraient ou
          infirmeraient le caractère fictif du fournisseur.
\end{enumerate}
\end{boxexo}

\begin{boxexo}[title={\ding{45}\quad Exercice~4.3. Un usage stratégique de NEK~? \niveauexpert}]
Reprenez la lecture de l'entrepôt de Douala (section~\ref{sec:c04-strategie}).
\begin{enumerate}[(a)]
    \item Recalculez $\mu$, $U$ et $H$ pour l'auto-évaluation et pour
          l'entretien formel.
    \item Proposez deux explications du déplacement des réponses, l'une
          sincère, l'autre stratégique.
    \item Concevez un protocole de recueil qui permettrait de départager ces
          deux explications.
    \item Rédigez la note que le chef de mission adresse au commanditaire, en
          respectant la conduite à tenir de la section~\ref{sec:c04-enquete}.
\end{enumerate}
\end{boxexo}

\bigskip

\begin{boxinfo}
\textbf{Transition.} La Partie~\ref{part:fondements} a posé ce que l'auditeur
doit savoir~: ce qu'est un audit, le contrôle interne, le risque et la
fraude. La Partie~\ref{part:methode} enseigne comment conduire la mission.
Le Chapitre~\ref{chap:mission} en présente les principes et le cycle.
\end{boxinfo}
